% !TeX root = ../../main.tex
% ============================================================
% Appendix A: Hierarchy（由 08_appendix.tex 拆分而来）
% 说明：本文件由 main.tex 合并进主文档，且位于 main.tex 的
%       \appendix 之后，因此 \section 会自动编号为 Appendix A。
%       内容忠实来自 策划案.docx 原文（英文）。
%       用 forest 横向树（根在左、向右展开）。
%       【自己加节点】想加子节点：在父节点 [ ] 里加一行 [名字]。
%       例如在 [EnemyPrefab (Example) 的 [ ] 里加 [NewBossAI]。
% ============================================================

% 附录起始页大标题（原在 main.tex，被 \include 的强制分页拆到单独一页；
% 移到本文件开头后与 Appendix A 标题同页）
\begin{center}
{\fontsize{16pt}{45pt}\selectfont\textbf{Appendix}}
\end{center}

\section{Hierarchy}   % 自动显示为 "Appendix A: Hierarchy"

\begin{center}
\begin{forest}
for tree={
  grow=east,
  draw,
  rounded corners=2pt,
  minimum height=2.1em,
  inner xsep=5pt,
  font=\footnotesize\sffamily,
  s sep=4pt,
  l sep=16pt,
  edge={gray, thick},
  parent anchor=east,
  child anchor=west,
  if n children=0{fill=gray!10}{fill=white}
}
[, phantom
  [EventSystem]
  [ParticleSystem]
  [PlayerUnit
    [SpriteRenderer]
    [PlayerAttack [Weapon]]
    [PlayerMovement]
    [PlayerHealth [Plug\_HealthBar]]
    [PlayerStamina [Plug\_StaminaBar]]
  ]
  [EnemyUnit
    [EnemyPrefab (Example)
      [SpriteRenderer]
      [Animation]
      [SpecialAl (If any)]
      [EnemyMovement]
      [EnemyAttack]
      [EnemyHealth]
      [DamagePopUp]
    ]
  ]
]
\end{forest}
\end{center}
